\documentclass[11pt]{article}
\usepackage[margin=1.1in]{geometry}
\usepackage{amsmath,amssymb,amsthm,booktabs,hyperref}
\hypersetup{colorlinks=true,linkcolor=blue,citecolor=blue,urlcolor=blue}

\newtheorem{theorem}{Theorem}
\newtheorem{lemma}[theorem]{Lemma}
\newtheorem{conjecture}[theorem]{Conjecture}
\theoremstyle{remark}
\newtheorem{remark}[theorem]{Remark}

\newcommand{\Jt}{\widetilde{J}}
\newcommand{\binomb}[2]{\tbinom{#1}{#2}}

\title{Off the Real Line\\[0.4em]\large A counterexample to the real-rootedness of Jack hit polynomials}
\author{Prajit Chennupati\\ \small\texttt{chennupati@gatech.edu} \and Aashrith Nalam\\ \small\texttt{analam3@gatech.edu}}
\date{November 24, 2026}



\begin{document}

\maketitle
\begin{abstract}
Alexandersson, Haglund and Wang conjectured that the coefficients $a_k(\mu,\lambda)$ in the expansion
$\langle \alpha^n J_\mu^{(1/\alpha)}, s_\lambda\rangle=\sum_k a_k(\mu,\lambda)\binom{\alpha+k}{n}$ are
nonnegative integers and that $\sum_k a_k(\mu,\lambda)z^k$ has only real zeros. We show that the
real-rootedness assertion fails for $\mu=(8,8)$, $\lambda=(2^8)$: the polynomial has degree $12$, a simple
zero at $0$, and exactly $9$ further real zeros, so two of its zeros are non-real. The proof is
computer-assisted and uses exact arithmetic only. The polynomial identities and the non-real-rootedness are
formally verified in Lean~4 with Mathlib. The positivity assertion holds for this pair, and an exhaustive
search finds no other failure with $n\le 17$.
\end{abstract}

\section{The conjecture}

Let $J_\mu^{(\alpha)}$ be the integral form Jack symmetric function, normalized so that the coefficient of
$m_{1^n}$ is $n!$ \cite{Mac,Sta}, and let $\langle\cdot,\cdot\rangle$ be the Hall inner product. Following
\cite{AHW18,AHW21} put
\[
  \Jt_\mu^{(\alpha)} := \alpha^{n} J_\mu^{(1/\alpha)}, \qquad n=|\mu| .
\]

\begin{conjecture}[{Alexandersson--Haglund--Wang \cite[Conj.~1]{AHW18,AHW21}}]\label{conj}
Let $\mu,\lambda$ be partitions of $n$ and write
$\langle \Jt_\mu^{(\alpha)}, s_\lambda\rangle=\sum_{k} a_k(\mu,\lambda)\binom{\alpha+k}{n}$. Then
\textup{(i)} $a_k(\mu,\lambda)\in\mathbb N$, and \textup{(ii)} the polynomial
$\sum_{k=0}^{n}a_k(\mu,\lambda)z^k$ has only real zeros.
\end{conjecture}

The polynomial in (ii) is called the \emph{Jack hit polynomial} on the SymCat database \cite{SymCat}.
Since $\{\binom{\alpha+k}{n}\}_{0\le k\le n}$ is a basis of polynomials of degree at most $n$, the $a_k$ are
well defined. (The source sums the expansion over $k\le n-1$ but the polynomial over $k\le n$; for the pair
below $a_0=a_{13}=\dots=a_{16}=0$, so this does not matter.)
Alexandersson, Haglund and Wang prove the case $\mu=(n)$ and the hook case $\mu=\lambda$, and report computer verification for
$n\le 11$ \cite{AHW18}. A companion Conjecture~2 uses the basis $\binom{\alpha}{k}k!$; part (ii) of
Conjecture~\ref{conj} implies part (ii) of Conjecture~2 \cite{AHW18}.

\section{The counterexample}

\begin{theorem}\label{thm}
Let $\mu=(8,8)$ and $\lambda=(2^8)$, so $n=16$. Then
\begin{align*}
p(\alpha):=\langle \Jt_\mu^{(\alpha)}, s_\lambda\rangle
 &= 112\,\alpha^2(\alpha-1)^2(\alpha-2)^2(\alpha-3)^2(\alpha+1)(2\alpha+1)\\
 &\qquad\times(104\alpha^6-684\alpha^5+2258\alpha^4-2901\alpha^3+8033\alpha^2+4530\alpha+12600),
\end{align*}
the coefficients $a_k=a_k(\mu,\lambda)$ are the nonnegative integers in Table~\ref{tab}, and the hit polynomial
$H(z)=\sum_{k=0}^{16}a_kz^k$ is not real-rooted. Precisely, $H(z)=z\,h(z)$ where $h$ has degree $11$, is
squarefree, and has exactly $9$ real zeros. Its non-real zeros are
\[
 -10.52898876935\ldots \pm 2.39070820941\ldots\, i .
\]
Hence part \textup{(ii)} of Conjecture~\ref{conj} is false.
\end{theorem}

\begin{table}[h]
\centering
\begin{tabular}{rr@{\qquad}rr}
\toprule
$k$ & $a_k$ & $k$ & $a_k$\\
\midrule
1 & 597\,793\,996\,800 & 7 & 146\,689\,299\,490\,406\,400\\
2 & 73\,236\,035\,174\,400 & 8 & 68\,260\,823\,569\,612\,800\\
3 & 2\,004\,440\,894\,668\,800 & 9 & 15\,222\,991\,343\,616\,000\\
4 & 19\,314\,619\,527\,168\,000 & 10 & 1\,568\,059\,637\,760\,000\\
5 & 79\,659\,352\,571\,904\,000 & 11 & 68\,516\,388\,864\,000\\
6 & 154\,554\,571\,254\,988\,800 & 12 & 804\,722\,688\,000\\
\bottomrule
\end{tabular}
\caption{The coefficients $a_k((8,8),(2^8))$; all other $a_k$ ($k=0$ and $13\le k\le16$) vanish. All
are divisible by $2\,090\,188\,800$.}\label{tab}
\end{table}

The real-rootedness failure has a short certificate, which is the part checked in Lean.

\begin{lemma}\label{lem}
If $f\in\mathbb R[x]$ is a real constant times a product of real linear factors $x-r_i$, then
$f'(x)^2-f(x)f''(x)\ge 0$ for all real $x$.
\end{lemma}

\begin{proof}
Write $Q_f(x)=f'(x)^2-f(x)f''(x)$. Constants have $Q=0$. If $f=(x-r)g$, then $f'=g+(x-r)g'$ and
$f''=2g'+(x-r)g''$, so $Q_f(x)=g(x)^2+(x-r)^2Q_g(x)$. Induct on the number of factors.
\end{proof}

For $H$ at $x=-10$ one computes exactly
\begin{align*}
H(-10)&=-83\,342\,017\,288\,983\,011\,328\,000,\\
H'(-10)&=68\,364\,158\,223\,859\,539\,148\,800,\\
H''(-10)&=-71\,278\,550\,656\,400\,528\,179\,200,
\end{align*}
so $Q_H(-10)=H'(-10)^2-H(-10)H''(-10)\approx-1.2668\times10^{45}<0$. By Lemma~\ref{lem}, $H$ does not have only
real zeros. The exact root count ($9$ real zeros of $h$) is a strengthening obtained by Sturm's theorem.

\section{Proof of Theorem~\ref{thm}}

The theorem combines two published facts with an exact computation.
\begin{itemize}
\item[(H1)] $P_\mu^{(\beta)}$ is the unique symmetric function of the form $m_\mu+\sum_{\rho<\mu}u_{\mu\rho}m_\rho$
(dominance order) orthogonal to all $P_\rho^{(\beta)}$, $\rho\neq\mu$, for
$\langle p_\sigma,p_\tau\rangle_\beta=\delta_{\sigma\tau}z_\sigma\beta^{\ell(\sigma)}$. Moreover
$J_\mu^{(\beta)}=c_\mu(\beta)P_\mu^{(\beta)}$ with $[m_{1^n}]J_\mu^{(\beta)}=n!$ \cite{Mac,Sta}.
\item[(H2)] The monomial coefficients of $J_\mu^{(\beta)}$ are polynomials in $\beta$ of degree at most $n$
\cite{KS97}. Hence $p(\alpha)$ is a polynomial of degree at most $16$.
\end{itemize}

\emph{Computation.} Fix $\beta=1/\alpha$ with $\alpha\in\{1,\dots,18\}$. The set $D=\{\rho:\rho<\mu\}$ is a
down-set, so $\mathrm{span}\{P_\rho:\rho\in D\}=\mathrm{span}\{m_\rho:\rho\in D\}$. Therefore the coefficients
$u_{\mu\rho}$ solve the linear system $\langle P_\mu,m_\rho\rangle_\beta=0$, $\rho\in D$. Its Gram matrix is
positive definite, so the solution is unique. We form the exact integer transition matrix between power sums and
monomials, solve the system over $\mathbb Q$, normalize by $[m_{1^n}]J=n!$, and compute
$\langle J,s_\lambda\rangle=\sum_\sigma c_\sigma\chi^\lambda(\sigma)$ for $J=\sum_\sigma c_\sigma p_\sigma$, with
characters from the Murnaghan--Nakayama rule. The $17$ values at $\alpha=1,\dots,17$ determine $p$ by (H2).
The value at $\alpha=18$ is a consistency check. Solving for the $a_k$ in the basis $\binom{\alpha+k}{16}$ gives
Table~\ref{tab}. Finally, $\gcd(h,h')=1$, Sturm's theorem gives $9$ real zeros of $h$, and FLINT/Arb rigorous
root isolation independently returns $9$ real zeros and the conjugate pair stated.

\emph{Independent confirmations.} Two further implementations reproduce the same $p(\alpha)$ and $a_k$
exactly. The first uses Stanley's Laplace--Beltrami recursion for the monomial coefficients of $P_\mu$, the hook
normalization, and the inverse Kostka matrix. The second was written independently: it takes
$J_\mu$ as the one-dimensional kernel of $D(\alpha)-e_\mu$ in the power-sum basis, uses dual Jacobi--Trudi for
$s_\lambda$, and interpolates at the disjoint points $\alpha=1009,\dots,1043$. The implementations reproduce the
proved cases $\langle\Jt_\mu,s_\mu\rangle=\prod_{s\in\mu}(\mathrm{arm}(s)+\alpha(\mathrm{leg}(s)+1))$ and the
$\mu=(n)$ formula of \cite{AHW18}, and the example $\mu=\lambda=(3)$ listed on \cite{SymCat}.

\section{Formal verification}

The Lean~4 file \texttt{lean-formalization/JackHit.lean} in the accompanying repository (see Reproducibility; Mathlib v4.36.0) proves the following.
\begin{itemize}
\item \texttt{p\_eq\_binomial\_expansion}: for every $\alpha\in\mathbb Q$,
$p(\alpha)=\sum_{k=0}^{16}a_k\binom{\alpha+k}{16}$ with the explicit $a_k$ of Table~\ref{tab}.
\item \texttt{hitPoly\_not\_realRooted}: $H$ is not of the form $c\prod_i(z-r_i)$ with $c,r_i\in\mathbb R$.
This goes through Lemma~\ref{lem} and the evaluation at $-10$.
\end{itemize}
Both theorems depend only on the axioms \texttt{propext}, \texttt{Classical.choice} and \texttt{Quot.sound}; no
\texttt{sorry} or \texttt{native\_decide} is used. The identity
$p(\alpha)=\langle\Jt_{(8,8)}^{(\alpha)},s_{(2^8)}\rangle$ is \emph{not} formalized, since Jack polynomials are
not available in Mathlib. It is the computer-verified hypothesis of the formal statement.

\section{Remarks}

\begin{remark}
An exhaustive exact search over all pairs $\mu,\lambda\vdash n$ with $n\le 17$ finds that parts (i) and (ii) of
Conjecture~\ref{conj} and of Conjecture~2 hold in every case except the pair of Theorem~\ref{thm}, which fails
only (ii) of Conjecture~\ref{conj}. Thus $n=16$ is the smallest size of a counterexample. This minimality
statement relies on one implementation only.
\end{remark}

\begin{remark}
In the family $\mu=(k,k)$, $\lambda=(2^k)$, the hit polynomial is real-rooted for $k\le 7$ and $k\in\{9,11\}$
and fails to be real-rooted for $k\in\{8,10,12\}$ (real zeros of the core versus its degree: $9/11$, $12/14$,
$15/17$). We do not know whether it fails for all even $k\ge 8$.
\end{remark}

\begin{remark}
Part (i) of Conjecture~\ref{conj} and both parts of Conjecture~2 (the ``rook'' version) remain open. The
SymCat remark that real-rootedness of hit polynomials would imply that of rook polynomials is consistent with
this example: the rook polynomial of this pair is real-rooted.
\end{remark}

\subsection*{Reproducibility}
{\sloppy All code and the Lean formalization are available at
\url{https://github.com/prajitchennupati/math/tree/main/jack-hit-polynomial-counterexample}; paths below are
relative to that folder. \texttt{paper/code/verify.py} recomputes Theorem~\ref{thm} from scratch in under one second (python-flint 0.9,
SymPy 1.14). \texttt{paper/code/jack.py} and \texttt{paper/code/search.py} implement the second method and the
exhaustive search, and \texttt{paper/code/independent/} contains the third implementation. The Lean
formalization is in \texttt{lean-formalization/}. This note has not been refereed.\par}

\begin{thebibliography}{9}
\bibitem{AHW18} P.~Alexandersson, J.~Haglund, G.~Wang, \emph{On the Schur expansion of Jack polynomials},
FPSAC 2018, S\'em. Lothar. Combin. 80B (2018), Art.~90; arXiv:1805.00511.
\bibitem{AHW21} P.~Alexandersson, J.~Haglund, G.~Wang, \emph{Some conjectures on the Schur expansion of Jack
polynomials}, J. Comb. 12 (2021), 215--233; arXiv:1810.12763.
\bibitem{KS97} F.~Knop, S.~Sahi, \emph{A recursion and a combinatorial formula for Jack polynomials},
Invent. Math. 128 (1997), 9--22.
\bibitem{Mac} I.~G.~Macdonald, \emph{Symmetric Functions and Hall Polynomials}, 2nd ed., Oxford, 1995.
\bibitem{Sta} R.~P.~Stanley, \emph{Some combinatorial properties of Jack symmetric functions},
Adv. Math. 77 (1989), 76--115.
\bibitem{SymCat} P.~Alexandersson, \emph{The symmetric functions catalog}, \url{https://www.symmetricfunctions.com/jack.htm}
(accessed November 2026).
\end{thebibliography}

\end{document}
